\section*{Bab \ref{ch:g7:oxygen-in-the-water} --- Oksigen di dalam Air}

\begin{solution}{exo:g7:oxygen-in-the-water:1}
Kira-kira \qty{10}{mg} per liter --- kasarnya tiga puluh kali lebih
sedikit daripada yang dibawa seliter udara.
\end{solution}

\begin{solution}{exo:g7:oxygen-in-the-water:2}
Udara di atasnya, yang melarut masuk di permukaannya --- paling cepat
di tempat airnya teraduk, terpercik dan berbuih; dan para \omterm{def:g5:food-webs:roles}{produsen} di
bawahnya, yang melepaskan \omterm{def:g7:what-is-respiration:respiration}{oksigen} pada siang hari --- paling cepat
dalam \omterm{def:g6:exploring-our-environment:conditions}{cahaya} yang terang dengan pertumbuhan hijau yang berlimpah.
\end{solution}

\begin{solution}{exo:g7:oxygen-in-the-water:3}
\omterm{def:g7:what-is-respiration:respiration}{Respirasi} \omterm{def:g1:animals-around-us:animal}{hewannya}, \omterm{def:g7:what-is-respiration:respiration}{respirasi} \omterm{def:g1:plants-around-us:plant}{tumbuhannya} sendiri pada malam hari, dan
\omterm{def:g7:what-is-respiration:respiration}{respirasi} para \omterm{def:g6:decomposers-and-soil:decomposers}{pengurainya} --- yang menjadi amat besar ketika materi
mati berlimpah.
\end{solution}

\begin{solution}{exo:g7:oxygen-in-the-water:4}
Air yang hangat memuat lebih sedikit \omterm{def:g7:oxygen-in-the-water:dissolved}{oksigen terlarut} (langit-langitnya
turun) sementara mesin setiap pelaku \omterm{def:g7:what-is-respiration:respiration}{respirasi} berjalan lebih cepat
(permintaannya naik): pasokan yang lebih sedikit, kebutuhan yang lebih
banyak, sekaligus.
\end{solution}

\begin{solution}{exo:g7:oxygen-in-the-water:5}
Pada siang hari pembuatan oleh \omterm{def:g5:food-webs:roles}{produsennya} menambahkan \omterm{def:g7:what-is-respiration:respiration}{oksigen} lebih
cepat daripada yang dikuras \omterm{def:g7:what-is-respiration:respiration}{respirasinya} --- lengkungnya menanjak ke
puncak sore harinya. Sepanjang malam setiap pelaku \omterm{def:g7:what-is-respiration:respiration}{respirasi} menimba
persediaannya sama sekali tanpa pembuatan --- lengkungnya meluncur ke
titik terendahnya tepat sebelum matahari terbit.
\end{solution}

\begin{solution}{exo:g7:oxygen-in-the-water:6}
Kaya: \omterm{prop:g3:sorting-living-things:groups}{ikan} trout (atau benih \omterm{prop:g3:sorting-living-things:groups}{ikan} trout) dan larva lalat sehari di
bawah batunya. Miskin: hanya \omterm{def:g5:biodiversity:species}{spesies} yang tahan --- \omterm{prop:g3:sorting-living-things:groups}{ikan} mas, siput
kolam, dan larva bersnorkel yang mengambil udara dari permukaannya.
\end{solution}

\begin{solution}{exo:g7:oxygen-in-the-water:7}
\omterm{def:g7:what-is-respiration:respiration}{Respirasi} malamnya --- \omterm{prop:g3:sorting-living-things:groups}{ikan}, \omterm{def:g1:plants-around-us:plant}{tumbuhan} dan \omterm{def:g6:decomposers-and-soil:decomposers}{pengurai} lumpurnya
bersama-sama --- menguras kolamnya sampai minimumnya tepat sebelum
matahari terbit, sehingga kolam yang terbebani mencekik \omterm{prop:g3:sorting-living-things:groups}{ikannya} pada
waktu fajar. Penawarnya: dari pemasok 1, aduklah permukaannya ---
jalankan sebuah air mancur atau air terjun buatan sepanjang malam;
dari pemasok 2, tipiskanlah \omterm{def:g1:plants-around-us:plant}{tumbuhan} yang terlalu rimbun dan
lumpurnya sehingga \omterm{def:g7:what-is-respiration:respiration}{respirasi} malamnya menyusut (dan produksi siang
harinya mencapai airnya, bukan hanya hamparan permukaannya).
\end{solution}

\begin{solution}{exo:g7:oxygen-in-the-water:8}
Limbahnya masuk --- para \omterm{def:g6:decomposers-and-soil:decomposers}{pengurainya} berpesta lalu berlipat ganda ---
\omterm{def:g7:what-is-respiration:respiration}{respirasi} massalnya menguras \omterm{def:g7:what-is-respiration:respiration}{oksigennya} --- \omterm{def:g5:biodiversity:species}{spesies} yang menuntut
lenyap dan yang tahan mengambil alih. Di hilirnya limbahnya akhirnya
dilahap, riamnya mengaerasi ulang airnya, dan \omterm{def:g5:biodiversity:species}{spesies} yang menuntut
itu kembali: pengurasnya tertutup, pemasoknya menyusul.
\end{solution}

\begin{solution}{exo:g7:oxygen-in-the-water:9}
Sebuah sensor mencatat satu saat; penghuninya mencatat sejarah. Sebuah
\omterm{def:g5:biodiversity:species}{spesies} yang menuntut hanya hidup di sebuah stasiun kalau airnya tetap
baik sepanjang seluruh masa tinggalnya --- termasuk pada fajar dan
minggu buruk yang tidak tertangkap satu kunjungan pun. Pendataannya
adalah minimum semusim, yang tertulis dalam bentuk \omterm{def:g1:animals-around-us:animal}{hewan}.
\end{solution}

\begin{solution}{exo:g7:oxygen-in-the-water:10}
Warna hijau seperti sup kacang itu adalah populasi besar \omterm{def:g5:first-look-at-cells:cell}{sel} tunggal
hijau: pada siang hari mereka menghasilkan, tetapi setiap malam yang
hangat seluruh supnya berespirasi --- dan minggu yang panas dan tenang
itu menurunkan langit-langitnya serta menghentikan pengisian ulang
permukaannya. Minimum fajarnya menuju nol: cuaca kematian \omterm{prop:g3:sorting-living-things:groups}{ikan}.
Jalankan air mancurnya paling keras sepanjang malam dan jam-jam
kecilnya --- siang hari dapat mengurus dirinya sendiri.
\end{solution}

\begin{solution}{exo:g7:oxygen-in-the-water:11}
\omterm{prop:g3:sorting-living-things:groups}{Ikan} trout adalah \omterm{def:g5:biodiversity:species}{spesies} yang menuntut: sungai yang dingin memuat
langit-langit tinggi yang dibutuhkannya; bendung dan saluran pacu yang
memercik menjaga pengisian ulangnya tetap tanpa henti melawan
pengurasan oleh \omterm{prop:g3:sorting-living-things:groups}{ikannya} yang berdesakan. Danau yang hangat dan tenang
gagal pada langit-langit dan pengisian ulangnya sekaligus --- negeri
\omterm{prop:g3:sorting-living-things:groups}{ikan} mas, bukan negeri \omterm{prop:g3:sorting-living-things:groups}{ikan} trout.
\end{solution}

\begin{solution}{exo:g7:oxygen-in-the-water:12}
Mata uang: setiap pelaku \omterm{def:g7:what-is-respiration:respiration}{respirasi} bawah air menimba satu persediaan
\omterm{def:g7:oxygen-in-the-water:dissolved}{terlarut} --- siapa tumbuh subur di mana dihargai dalam \omterm{def:g7:what-is-respiration:respiration}{oksigen}
(\cref{prop:g7:oxygen-in-the-water:map}). Jumlah uang yang beredar:
\omterm{def:g6:exploring-our-environment:conditions}{suhu} menetapkan berapa banyak yang dapat dimuat airnya --- dingin
menaikkan langit-langitnya, kehangatan menurunkannya
(\cref{prop:g7:oxygen-in-the-water:drains}). Pembelanja terbesar:
kalau ada materi mati, \omterm{def:g7:what-is-respiration:respiration}{respirasi} massal para \omterm{def:g6:decomposers-and-soil:decomposers}{pengurainya} menguras
melebihi siapa pun --- itulah pelajaran dari pipa pabrik olahan
\omterm{def:g2:animals-grow-up:born}{susunya} (\cref{ex:g7:oxygen-in-the-water:waste}).
\end{solution}

\begin{solution}{pb:g7:oxygen-in-the-water:1}
\textbf{1.} K1 dingin --- langit-langit yang tinggi --- dan teraduk
oleh bendungnya --- pengisian ulang permukaan pada kecepatan penuh:
bacaan yang dekat langit-langitnya. K2 hangat --- langit-langit yang
turun --- dan tidak teraduk: sebuah bacaan yang sederhana adalah semua
yang dapat dipegang pemasoknya.

\textbf{2.} \omterm{def:g5:biodiversity:species}{Spesies} yang menuntut di K1 menjamin \omterm{def:g7:what-is-respiration:respiration}{oksigen} yang kaya dan
mantap; rombongan yang tahan di K2 --- \omterm{prop:g3:sorting-living-things:groups}{ikan} mas, siput, para pemakai
snorkel --- menandai sebuah perairan yang berjalan miskin, setidaknya
pada jam-jam buruknya.

\textbf{3.} \omterm{def:g7:what-is-respiration:respiration}{Respirasi} para \omterm{def:g6:decomposers-and-soil:decomposers}{pengurainya}: air buangan pabrik olahan
\omterm{def:g2:animals-grow-up:born}{susunya} yang kaya organik memberi makan sebuah populasi massal
mereka, dan pesta mereka menguras persediaannya jauh di bawah apa yang
dapat dijelaskan \omterm{def:g6:exploring-our-environment:conditions}{suhunya} sendirian.

\textbf{4.} Dua kilometer pelahapan sudah menghabiskan limbahnya ---
pestanya usai dan para pembelanjanya menipis --- sementara riam
berbatunya sudah mengaduk \omterm{def:g7:what-is-respiration:respiration}{oksigen} kembali masuk: kedua pemasoknya
pulih, pengurasnya tertutup, dan pendataannya menjamin pemulihannya.

\textbf{5.} Sebuah gelombang harian: palung fajar pada \qty{4}{mg/L},
puncak sore pada \qty{9}{mg/L}. Kedua pedagangnya: pembuatan siang
hari (yang menambah pada siang hari) dan \omterm{def:g7:what-is-respiration:respiration}{respirasi} sepanjang waktu
(yang menguras sendirian sepanjang malam).

\textbf{6.} Dinginnya musim dingin menaikkan langit-langitnya,
sehingga airnya memuat lebih banyak; dan dingin memperlambat setiap
pelaku \omterm{def:g7:what-is-respiration:respiration}{respirasi} --- \omterm{prop:g3:sorting-living-things:groups}{ikan}, \omterm{def:g1:plants-around-us:plant}{tumbuhan}, \omterm{def:g6:decomposers-and-soil:decomposers}{pengurai} lumpurnya --- sehingga
pengurasnya menyusut. Lebih mantap pula, karena gelombang harian
\omterm{def:g5:food-webs:roles}{produsennya} lebih kecil.

\textbf{7.} K3 lebih dahulu (sudah terkuras, dan panasnya mempercepat
\omterm{def:g6:decomposers-and-soil:decomposers}{pengurainya}); K2 kedua (langit-langit yang rendah, air yang tenang,
palung fajarnya); K4 ketiga (sedang pulih tetapi berada di hilir
masalahnya); K1 terakhir (dingin dan teraduk --- bendungnya terus
mengisi ulang apa pun cuacanya).

\textbf{8.} Sebelum matahari terbit adalah minimum hariannya: bacaan
pada waktu itu mencatat keadaan terburuk yang harus dilalui kehidupan
sungainya. Jadwal sore akan mengarsipkan tahunnya menurut saat-saat
terbaiknya.

\textbf{9.} Sensornya pulih dalam hitungan minggu --- pengurasnya
tertutup begitu makanan bagi \omterm{def:g6:decomposers-and-soil:decomposers}{pengurainya} berhenti. Pendataannya pulih
lebih lambat dan berurutan: mula-mula rombongan yang tahan menipis,
lalu para penjajah yang hanyut dan terbang
(logika \cref{ch:g6:how-plants-spread}, versi \omterm{def:g1:animals-around-us:animal}{hewan}) menyemai ulang
\omterm{def:g5:biodiversity:species}{spesies} yang menuntut --- larva lalat sehari lebih dahulu dari
hulunya, \omterm{prop:g3:sorting-living-things:groups}{ikan} trout terakhir.

\textbf{10.} Naungannya juga menjaga ruasnya tetap sejuk --- dan
dinginlah langit-langitnya. Penebangannya akan menambah produksi siang
hari tetapi menurunkan langit-langitnya serta mempercepat setiap
pengurasnya di dalam air yang menghangat; pada sebuah ruas yang lambat
kerugiannya mungkin melebihi keuntungannya. Jawabannya: pertahankan
pohonnya, aduklah airnya sebagai gantinya.

\textbf{11.} Pendataannya: larva lalat sehari tidak dapat memalsukan
sebuah musim yang baik --- mereka menjalani setiap fajarnya. Bacaan
sore yang biasa-biasa saja itu satu saat, barangkali sebuah hari yang
sedang tidak baik; larvanya adalah catatan bersumpah tentang musimnya.

\textbf{12.} ``Jagalah sungainya tetap dingin, teraduk dan tidak
terbebani, maka \omterm{def:g7:what-is-respiration:respiration}{oksigennya} --- dan segala yang menghirupnya --- akan
mengurus dirinya sendiri.''
\end{solution}
